\chapter{Brakes}
% full version: chapters/03_gaba.tex

Let us add the second most numerous type of neuron to the machine: \nt{GABA}. Its signal does not nudge the neighbor toward a click; it pushes it away. These are the brakes. There are not many of them, from a fifth to a quarter of all neurons in the cortex, but without them the machine does not run.

There is one limitation: a ``plus'' neuron cannot brake a neighbor by itself. If you need a negative connection, you have to build it through a go-between: the ``plus'' wakes up a braking neuron, and that one brakes the target. Flipping the sign costs one neuron and a small delay.

\section*{``But Not If''}

The most useful operation with brakes is a veto: ``fire on A, but not if B has come.'' The neuron gets a push from A and a brake from a go-between that B wakes up. Vetoes and ``or'' are enough to build any logic.

\pencilfig{3}{\pencilart{3}}{``A, but not if B'': the red neuron blocks the blue one's path.}

A veto in time gives the motion detector that sits in your eye. Two neighboring sensors: one excites the output neuron, the other inhibits it with a delay. If a spot of light runs in one direction at the right speed, the inhibition arrives exactly when the excitation does and cancels it. If it runs the other way, the inhibition is late, and the neuron manages to click. The result is a neuron that sees motion in one direction only.

\pencilfig{103}{%
% sensors on top; the left one excites the output, the right one inhibits it through a GABA go-between and a delay
\draw[dotpen=pcAmber, wash=pcAmber] (0,1.6) circle (0.16); \draw[dotpen=pcAmber, wash=pcAmber] (1.6,1.6) circle (0.16);
\node[lbl, anchor=south] at (0.8,1.8) {two sensors};
\draw[pen=pcBlue, wash=pcBlue] (0.4,0) circle (0.24);
\draw[arr=pcBlue] (0,1.44) -- (0.3,0.25);
\draw[pen=pcBlue] (1.6,1.44) -- (1.6,1.18);
\draw[penlite, fill=white] (0.95,0.9) rectangle (2.25,1.18); \node[font=\tiny\itshape, text=pcGray] at (1.6,1.04) {delay};
\draw[pen=pcBlue] (1.6,0.9) -- (1.6,0.72);
\draw[pen=pcRed, wash=pcRed] (1.6,0.55) circle (0.17);
\draw[pcRed, line width=0.7pt, -{Bar[width=5pt]}] (1.45,0.45) -- (0.66,0.08);
\draw[arr] (0.4,-0.24) -- (0.4,-0.6); \node[lbl, anchor=north] at (0.4,-0.6) {output};
% outcomes
\begin{scope}[xshift=3.1cm]
  \draw[dotpen=pcAmber, wash=pcAmber] (0,1.25) circle (0.1); \draw[arr=pcAmber] (0.18,1.25) -- (1.1,1.25);
  \node[lbl, anchor=west] at (1.25,1.25) {left to right: click};
  \draw[dotpen=pcAmber, wash=pcAmber] (1.1,0.45) circle (0.1); \draw[arr=pcAmber] (0.92,0.45) -- (0,0.45);
  \node[lbl, anchor=west] at (1.25,0.45) {right to left: silence};
  \node[lbl, anchor=west, align=left] at (0,-0.35) {the inhibition arrives\\just with the excitation};
\end{scope}
}{A motion detector: the inhibition arrives with a delay and cancels the response only if the spot runs from right to left.}


\section*{Subtraction or Division}

Inhibition can do more than subtract. An open braking channel does not pull the voltage down; it pulls it \emph{toward rest}, as if a second hole had been punched in the bucket. Because of this, any inflow of water raises the level less. Inhibition does not take a fixed portion away from the signal; it \emph{divides} it: this is a volume control. If the strength of inhibition grows with the overall activity around, each neuron responds not to how strong its own signal is but to how strong it is \emph{compared with its neighbors}. That is why the same network works both in dim light and in bright sunshine. True, such a ``divider'' acts cleanly on the spike rate only together with steady background noise, which is always present in a living brain.

\pencilfig{104}{%
\foreach \dx/\t in {0/subtraction,4.6/division}{
  \begin{scope}[xshift=\dx cm]
  \draw[pen] (0,0) -- (3.6,0); \draw[pen] (0,0) -- (0,1.9);
  \node[lbl, anchor=north] at (1.8,-0.05) {signal}; \node[lbl, anchor=south, rotate=90] at (-0.05,0.95) {response};
  \draw[pen=pcBlue] (0.4,0) -- (3.3,1.75);
  \node[font=\small\itshape, text=pcGray, anchor=south] at (1.8,1.95) {\t};
  \end{scope}}
\draw[pen=pcRed, line width=0.9pt] (1.3,0) -- (3.4,1.27);
\node[lbl, text=pcRed, anchor=west] at (2.25,0.35) {shift};
\begin{scope}[xshift=4.6cm]
\draw[pen=pcRed, line width=0.9pt] (0.4,0) -- (3.3,0.8);
\node[lbl, text=pcRed, anchor=west] at (2.0,0.25) {lower slope};
\end{scope}
\node[lbl, text=pcBlue, anchor=east] at (2.25,1.45) {no brake};
}{Blue is the neuron's response without inhibition, red is with it. Subtraction shifts the threshold; division turns down the volume, like a control knob.}

There is a second consequence: the extra hole makes the bucket forget faster. The window within which two signals must land to add up gets narrower. Inhibition makes coincidence detectors stricter.

\section*{Choice}

Now for the main thing that was missing: choosing one out of many. Two groups of neurons inhibit each other. While the inhibition is weak, both work, just more quietly. But if it is strong enough, any small difference grows, as on a seesaw, until one group falls completely silent. Winner takes all. And the win holds: for the network to change its mind, a new argument must be noticeably stronger than the old one. The machine does not jump at every rustle.

Brakes also erase memory: on a ``reset'' signal a braking neuron puts out the campfire from the last chapter. Two such cells that inhibit each other are the counterpart of a latch, the basis of all computer memory.

\section*{Life on the Edge}

A network of ``pluses'' alone can run away into an avalanche: each spike gives rise to more than one next spike. The brakes hold it at its operating point, the way a thermostat holds the temperature of an oven. There are two conditions: the inhibition must be strong enough and fast enough. If it is strong but slow, the regulator starts to overshoot, like someone twisting the shower tap without waiting for the hot water to arrive. The temperature swings back and forth. In the brain such a loop of ``pluses'' and ``minuses'' produces a fast rhythm, several dozen oscillations per second.

The cortex, it seems, lives right on the edge: its own feedback is strong enough to run away into an avalanche without the brakes. Only the fast brakes hold it back, and that is why it responds so sensitively to weak signals.

The chapter's conclusion is unexpected: brakes add almost no new computations to the machine. They add \emph{control}: choice, reset, volume control, stability. Excitation carries the data; inhibition decides which data get through, when, and how loudly.

\fullversion{3}
